\documentclass[11pt,a4paper]{article}
\usepackage[utf8]{inputenc}
\usepackage[T1]{fontenc}
\usepackage[spanish,es-minimal,es-noshorthands]{babel}
\usepackage{lmodern}
\usepackage[sfdefault,lining]{FiraSans}
\usepackage[a4paper,top=2.3cm,bottom=2.2cm,left=1.8cm,right=1.8cm,headheight=1.3cm,headsep=0.4cm,footskip=1cm]{geometry}
\usepackage{amsmath,amssymb}
\allowdisplaybreaks
\usepackage{icomma}
\usepackage{graphicx}
\usepackage[table]{xcolor}
\usepackage{tikz}
\usepackage{fontawesome5}
\usepackage{ccicons}
\usepackage{fancyhdr}
\usepackage{lastpage}
\usepackage{tasks}
\usepackage{needspace}
\usepackage[most]{tcolorbox}
\usepackage{array}
\usepackage[hidelinks]{hyperref}
\AtBeginDocument{\hypersetup{pdftitle={Soluciones de la ficha de la unidad 6: Cálculo de derivadas. Derivabilidad de una función. 2º Bach.},pdfsubject={Matemáticas II 2º Bach.},pdfkeywords={derivada, recta tangente, recta normal, regla de la cadena, derivación logarítmica, derivación implícita, derivabilidad, funciones a trozos, parámetros, PAU}}}

\definecolor{navy}{HTML}{1B2748}
\definecolor{accent}{HTML}{E8603A}
\definecolor{teal}{HTML}{1F9E90}
\definecolor{soft}{HTML}{EEF1F8}
\definecolor{muted}{HTML}{5D6685}

\setlength{\parindent}{0pt}
\setlength{\parskip}{3pt}
\settasks{label=\alph*), label-format=\bfseries\color{navy}, label-width=1.3em, item-indent=1.8em, column-sep=0.6em, after-item-skip=5pt, before-skip=3pt, after-skip=4pt}

% ---------- capçalera i peu ----------
\newcommand{\llicencia}{\href{https://creativecommons.org/licenses/by-nc-sa/4.0/deed.es}{\ccbyncsa}\ \ Licencia Creative Commons Atribución-NoComercial-CompartirIgual 4.0 Internacional (CC BY-NC-SA 4.0)}
\pagestyle{fancy}
\fancyhf{}
\fancyhead[L]{\raisebox{-0.35\height}{\includegraphics[height=0.95cm]{logo-gva.png}}}
\fancyhead[C]{\small\color{muted}Soluciones de la ficha · Unidad 6: Cálculo de derivadas · 2\textsuperscript{o} Bach.}
\fancyhead[R]{\raisebox{-0.35\height}{\includegraphics[height=0.95cm]{logo-vermellar.png}}}
\fancyfoot[L]{\scriptsize\color{muted}\llicencia}
\fancyfoot[R]{\small\color{muted}\thepage}
\renewcommand{\headrulewidth}{0.4pt}
\renewcommand{\headrule}{\color{muted!40}\hrule height \headrulewidth}
\fancypagestyle{primera}{\fancyhf{}\renewcommand{\headrulewidth}{0pt}
  \fancyfoot[L]{\scriptsize\color{muted}\llicencia}
  \fancyfoot[R]{\small\color{muted}\thepage}}

% ---------- nivells i exercicis ----------
\newcommand{\estrelles}[1]{{\color{accent}\ifcase#1\or\faStar\faStar[regular]\faStar[regular]\or\faStar\faStar\faStar[regular]\or\faStar\faStar\faStar\fi}}
\newcounter{exer}
\newcommand{\ex}[1]{\par\needspace{4\baselineskip}\medskip\refstepcounter{exer}\noindent{\color{navy}\bfseries\theexer.}\ \ }

\newcommand{\seccio}[2]{\par\needspace{10\baselineskip}\bigskip
  \begin{tcolorbox}[colback=navy,colframe=navy,coltext=white,arc=2mm,boxrule=0pt,left=3mm,right=3mm,top=1.2mm,bottom=1.2mm]
  {\large\bfseries #1\quad #2}
  \end{tcolorbox}\vspace{-2pt}}


\setlength{\emergencystretch}{3em}

% quadre de recordatori (només a la fitxa d'exercicis)
\newcommand{\recorda}[1]{\par\smallskip\begin{tcolorbox}[colback=accent!7,colframe=accent!55,arc=2mm,boxrule=0.5pt,left=2.5mm,right=2.5mm,top=1mm,bottom=1mm]\small{\color{accent}\faLightbulb[regular]}\ \textbf{Recuerda:}\ #1\end{tcolorbox}}
% etiqueta dels exercicis de la PAU d'altres comunitats
\newcommand{\pau}[1]{{\color{teal}\small\bfseries\faUniversity\ PAU · #1.}\ }
\newcommand{\res}[1]{\boxed{#1}}\newcommand{\rest}[1]{\fbox{#1}}

% =====================================================================
% Capçalera FentMates de les fitxes: un color per a cada curs.
% Els colors són els de la pàgina inicial (LEVELS a assets/js/materials.js);
% el logo ve de _eines/logo/pdf/ (python3 _eines/logo/genera.py).
%
% Ús, al final del preàmbul de la unitat:
%   % =====================================================================
% Capçalera FentMates de les fitxes: un color per a cada curs.
% Els colors són els de la pàgina inicial (LEVELS a assets/js/materials.js);
% el logo ve de _eines/logo/pdf/ (python3 _eines/logo/genera.py).
%
% Ús, al final del preàmbul de la unitat:
%   % =====================================================================
% Capçalera FentMates de les fitxes: un color per a cada curs.
% Els colors són els de la pàgina inicial (LEVELS a assets/js/materials.js);
% el logo ve de _eines/logo/pdf/ (python3 _eines/logo/genera.py).
%
% Ús, al final del preàmbul de la unitat:
%   \input{../comu/fentmates.tex}
%   \fmcurs{1eso}
% i, al començament del document:
%   \capcalera{Matemàtiques}{1r ESO}{Fitxa d'exercicis · Unitat 2}{Divisibilitat}
% Opcional abans de \input: \def\fmneed{14} (espai mínim abans de cada secció).
% =====================================================================
\makeatletter
\def\fm@logodir{../../logo/pdf}
\providecommand{\fmneed}{10}
\def\fm@cursnom{}

% clau del curs, color i color del text damunt del color
\newcommand{\fm@curs}[3]{\@namedef{fm@color@#1}{#2}\@namedef{fm@text@#1}{#3}}
\fm@curs{1eso}{E6B84A}{navy}
\fm@curs{2eso}{3B82F6}{white}
\fm@curs{3eso}{FF7A50}{white}
\fm@curs{4eso-a}{A78BFA}{white}
\fm@curs{4eso-b}{6D28D9}{white}
\fm@curs{1batx-mi}{9B1C1C}{white}
\fm@curs{1batx-ccss}{EF6B6B}{white}
\fm@curs{2batx-mii}{1E6B47}{white}
\fm@curs{2batx-ccss}{2BB3A3}{white}

\newcommand{\fmcurs}[1]{%
  \@ifundefined{fm@color@#1}{\PackageError{fentmates}{Curs desconegut: #1}{}}{}%
  \edef\fm@tmp{\noexpand\definecolor{curs}{HTML}{\@nameuse{fm@color@#1}}%
               \noexpand\colorlet{curstext}{\@nameuse{fm@text@#1}}}\fm@tmp
  \colorlet{cursfosc}{curs!72!black}%
  \def\fm@logo{\fm@logodir/vertical-#1.pdf}%
  \def\fm@marca{\fm@logodir/marca-#1.pdf}}

% ---------- capçalera de la primera pàgina ----------
% #1 matèria (petit), #2 curs, #3 tipus de document i unitat, #4 títol
\newcommand{\capcalera}[4]{%
  \gdef\fm@cursnom{#2}%
  \par\noindent
  \begin{tcolorbox}[colback=curs,colframe=curs,boxrule=0pt,arc=2.2mm,width=3.5cm,height=2.15cm,
      valign=center,left=3mm,right=2mm,top=1mm,bottom=1mm,nobeforeafter,box align=center]
    \color{curstext}{\footnotesize #1}\par\vspace{1pt}{\fontsize{19}{21}\selectfont\bfseries #2\par}%
  \end{tcolorbox}\hspace{5mm}%
  \begin{minipage}[c]{\dimexpr\linewidth-3.5cm-5mm-2.9cm\relax}
    {\color{cursfosc}\small\bfseries\MakeUppercase{#3}}\par\vspace{2pt}
    {\color{navy}\fontsize{20}{23}\selectfont\bfseries #4\par}
  \end{minipage}\hfill
  \begin{minipage}[c]{2.6cm}\raggedleft\includegraphics[width=2.6cm]{\fm@logo}\end{minipage}\par
  \vspace{2pt}}

% ---------- barres de secció del color del curs ----------
\renewcommand{\seccio}[2]{\par\needspace{\fmneed\baselineskip}\bigskip
  \begin{tcolorbox}[colback=curs,colframe=curs,coltext=curstext,arc=2mm,boxrule=0pt,left=3mm,right=3mm,top=1.2mm,bottom=1.2mm]
  {\large\bfseries #1\quad #2}
  \end{tcolorbox}\vspace{-2pt}}

% ---------- estrelles de dificultat del color del curs (a Batxillerat no se'n posen) ----------
\@ifundefined{estrelles}{}{%
  \renewcommand{\estrelles}[1]{{\color{curs}\ifcase#1\or\faStar\faStar[regular]\faStar[regular]\or\faStar\faStar\faStar[regular]\or\faStar\faStar\faStar\fi}}}

% ---------- capçalera i peu de les pàgines ----------
% primera pàgina: logos institucionals al peu
\fancypagestyle{primera}{\fancyhf{}\renewcommand{\headrulewidth}{0pt}
  \fancyfoot[L]{\raisebox{-0.32\height}{\includegraphics[height=0.78cm]{logo-gva.png}}\hspace{5mm}%
    \raisebox{-0.32\height}{\includegraphics[height=0.78cm]{logo-vermellar.png}}}
  \fancyfoot[R]{\scriptsize\color{muted}\href{https://creativecommons.org/licenses/by-nc-sa/4.0/deed.ca}{\ccbyncsa}\ \ CC BY-NC-SA 4.0 · fentmates.org\qquad\small\thepage}}
% la resta: el símbol a l'esquerra i el curs a la dreta, amb el títol al centre com abans
\fancyhead[L]{\raisebox{-0.3\height}{\includegraphics[height=0.62cm]{\fm@marca}}}
\fancyhead[R]{\small\bfseries\color{cursfosc}\fm@cursnom}
\renewcommand{\headrule}{\color{curs!45}\hrule height \headrulewidth}
\makeatother

%   \fmcurs{1eso}
% i, al començament del document:
%   \capcalera{Matemàtiques}{1r ESO}{Fitxa d'exercicis · Unitat 2}{Divisibilitat}
% Opcional abans de \input: \def\fmneed{14} (espai mínim abans de cada secció).
% =====================================================================
\makeatletter
\def\fm@logodir{../../logo/pdf}
\providecommand{\fmneed}{10}
\def\fm@cursnom{}

% clau del curs, color i color del text damunt del color
\newcommand{\fm@curs}[3]{\@namedef{fm@color@#1}{#2}\@namedef{fm@text@#1}{#3}}
\fm@curs{1eso}{E6B84A}{navy}
\fm@curs{2eso}{3B82F6}{white}
\fm@curs{3eso}{FF7A50}{white}
\fm@curs{4eso-a}{A78BFA}{white}
\fm@curs{4eso-b}{6D28D9}{white}
\fm@curs{1batx-mi}{9B1C1C}{white}
\fm@curs{1batx-ccss}{EF6B6B}{white}
\fm@curs{2batx-mii}{1E6B47}{white}
\fm@curs{2batx-ccss}{2BB3A3}{white}

\newcommand{\fmcurs}[1]{%
  \@ifundefined{fm@color@#1}{\PackageError{fentmates}{Curs desconegut: #1}{}}{}%
  \edef\fm@tmp{\noexpand\definecolor{curs}{HTML}{\@nameuse{fm@color@#1}}%
               \noexpand\colorlet{curstext}{\@nameuse{fm@text@#1}}}\fm@tmp
  \colorlet{cursfosc}{curs!72!black}%
  \def\fm@logo{\fm@logodir/vertical-#1.pdf}%
  \def\fm@marca{\fm@logodir/marca-#1.pdf}}

% ---------- capçalera de la primera pàgina ----------
% #1 matèria (petit), #2 curs, #3 tipus de document i unitat, #4 títol
\newcommand{\capcalera}[4]{%
  \gdef\fm@cursnom{#2}%
  \par\noindent
  \begin{tcolorbox}[colback=curs,colframe=curs,boxrule=0pt,arc=2.2mm,width=3.5cm,height=2.15cm,
      valign=center,left=3mm,right=2mm,top=1mm,bottom=1mm,nobeforeafter,box align=center]
    \color{curstext}{\footnotesize #1}\par\vspace{1pt}{\fontsize{19}{21}\selectfont\bfseries #2\par}%
  \end{tcolorbox}\hspace{5mm}%
  \begin{minipage}[c]{\dimexpr\linewidth-3.5cm-5mm-2.9cm\relax}
    {\color{cursfosc}\small\bfseries\MakeUppercase{#3}}\par\vspace{2pt}
    {\color{navy}\fontsize{20}{23}\selectfont\bfseries #4\par}
  \end{minipage}\hfill
  \begin{minipage}[c]{2.6cm}\raggedleft\includegraphics[width=2.6cm]{\fm@logo}\end{minipage}\par
  \vspace{2pt}}

% ---------- barres de secció del color del curs ----------
\renewcommand{\seccio}[2]{\par\needspace{\fmneed\baselineskip}\bigskip
  \begin{tcolorbox}[colback=curs,colframe=curs,coltext=curstext,arc=2mm,boxrule=0pt,left=3mm,right=3mm,top=1.2mm,bottom=1.2mm]
  {\large\bfseries #1\quad #2}
  \end{tcolorbox}\vspace{-2pt}}

% ---------- estrelles de dificultat del color del curs (a Batxillerat no se'n posen) ----------
\@ifundefined{estrelles}{}{%
  \renewcommand{\estrelles}[1]{{\color{curs}\ifcase#1\or\faStar\faStar[regular]\faStar[regular]\or\faStar\faStar\faStar[regular]\or\faStar\faStar\faStar\fi}}}

% ---------- capçalera i peu de les pàgines ----------
% primera pàgina: logos institucionals al peu
\fancypagestyle{primera}{\fancyhf{}\renewcommand{\headrulewidth}{0pt}
  \fancyfoot[L]{\raisebox{-0.32\height}{\includegraphics[height=0.78cm]{logo-gva.png}}\hspace{5mm}%
    \raisebox{-0.32\height}{\includegraphics[height=0.78cm]{logo-vermellar.png}}}
  \fancyfoot[R]{\scriptsize\color{muted}\href{https://creativecommons.org/licenses/by-nc-sa/4.0/deed.ca}{\ccbyncsa}\ \ CC BY-NC-SA 4.0 · fentmates.org\qquad\small\thepage}}
% la resta: el símbol a l'esquerra i el curs a la dreta, amb el títol al centre com abans
\fancyhead[L]{\raisebox{-0.3\height}{\includegraphics[height=0.62cm]{\fm@marca}}}
\fancyhead[R]{\small\bfseries\color{cursfosc}\fm@cursnom}
\renewcommand{\headrule}{\color{curs!45}\hrule height \headrulewidth}
\makeatother

%   \fmcurs{1eso}
% i, al començament del document:
%   \capcalera{Matemàtiques}{1r ESO}{Fitxa d'exercicis · Unitat 2}{Divisibilitat}
% Opcional abans de \input: \def\fmneed{14} (espai mínim abans de cada secció).
% =====================================================================
\makeatletter
\def\fm@logodir{../../logo/pdf}
\providecommand{\fmneed}{10}
\def\fm@cursnom{}

% clau del curs, color i color del text damunt del color
\newcommand{\fm@curs}[3]{\@namedef{fm@color@#1}{#2}\@namedef{fm@text@#1}{#3}}
\fm@curs{1eso}{E6B84A}{navy}
\fm@curs{2eso}{3B82F6}{white}
\fm@curs{3eso}{FF7A50}{white}
\fm@curs{4eso-a}{A78BFA}{white}
\fm@curs{4eso-b}{6D28D9}{white}
\fm@curs{1batx-mi}{9B1C1C}{white}
\fm@curs{1batx-ccss}{EF6B6B}{white}
\fm@curs{2batx-mii}{1E6B47}{white}
\fm@curs{2batx-ccss}{2BB3A3}{white}

\newcommand{\fmcurs}[1]{%
  \@ifundefined{fm@color@#1}{\PackageError{fentmates}{Curs desconegut: #1}{}}{}%
  \edef\fm@tmp{\noexpand\definecolor{curs}{HTML}{\@nameuse{fm@color@#1}}%
               \noexpand\colorlet{curstext}{\@nameuse{fm@text@#1}}}\fm@tmp
  \colorlet{cursfosc}{curs!72!black}%
  \def\fm@logo{\fm@logodir/vertical-#1.pdf}%
  \def\fm@marca{\fm@logodir/marca-#1.pdf}}

% ---------- capçalera de la primera pàgina ----------
% #1 matèria (petit), #2 curs, #3 tipus de document i unitat, #4 títol
\newcommand{\capcalera}[4]{%
  \gdef\fm@cursnom{#2}%
  \par\noindent
  \begin{tcolorbox}[colback=curs,colframe=curs,boxrule=0pt,arc=2.2mm,width=3.5cm,height=2.15cm,
      valign=center,left=3mm,right=2mm,top=1mm,bottom=1mm,nobeforeafter,box align=center]
    \color{curstext}{\footnotesize #1}\par\vspace{1pt}{\fontsize{19}{21}\selectfont\bfseries #2\par}%
  \end{tcolorbox}\hspace{5mm}%
  \begin{minipage}[c]{\dimexpr\linewidth-3.5cm-5mm-2.9cm\relax}
    {\color{cursfosc}\small\bfseries\MakeUppercase{#3}}\par\vspace{2pt}
    {\color{navy}\fontsize{20}{23}\selectfont\bfseries #4\par}
  \end{minipage}\hfill
  \begin{minipage}[c]{2.6cm}\raggedleft\includegraphics[width=2.6cm]{\fm@logo}\end{minipage}\par
  \vspace{2pt}}

% ---------- barres de secció del color del curs ----------
\renewcommand{\seccio}[2]{\par\needspace{\fmneed\baselineskip}\bigskip
  \begin{tcolorbox}[colback=curs,colframe=curs,coltext=curstext,arc=2mm,boxrule=0pt,left=3mm,right=3mm,top=1.2mm,bottom=1.2mm]
  {\large\bfseries #1\quad #2}
  \end{tcolorbox}\vspace{-2pt}}

% ---------- estrelles de dificultat del color del curs (a Batxillerat no se'n posen) ----------
\@ifundefined{estrelles}{}{%
  \renewcommand{\estrelles}[1]{{\color{curs}\ifcase#1\or\faStar\faStar[regular]\faStar[regular]\or\faStar\faStar\faStar[regular]\or\faStar\faStar\faStar\fi}}}

% ---------- capçalera i peu de les pàgines ----------
% primera pàgina: logos institucionals al peu
\fancypagestyle{primera}{\fancyhf{}\renewcommand{\headrulewidth}{0pt}
  \fancyfoot[L]{\raisebox{-0.32\height}{\includegraphics[height=0.78cm]{logo-gva.png}}\hspace{5mm}%
    \raisebox{-0.32\height}{\includegraphics[height=0.78cm]{logo-vermellar.png}}}
  \fancyfoot[R]{\scriptsize\color{muted}\href{https://creativecommons.org/licenses/by-nc-sa/4.0/deed.ca}{\ccbyncsa}\ \ CC BY-NC-SA 4.0 · fentmates.org\qquad\small\thepage}}
% la resta: el símbol a l'esquerra i el curs a la dreta, amb el títol al centre com abans
\fancyhead[L]{\raisebox{-0.3\height}{\includegraphics[height=0.62cm]{\fm@marca}}}
\fancyhead[R]{\small\bfseries\color{cursfosc}\fm@cursnom}
\renewcommand{\headrule}{\color{curs!45}\hrule height \headrulewidth}
\makeatother

\fmcurs{2batx-mii}

\begin{document}
\thispagestyle{primera}

\capcalera{Matemáticas II}{2\textsuperscript{o} Bach.}{Soluciones de la ficha · Unidad 6}{\raggedright Cálculo de derivadas. Derivabilidad de una función}
\vspace{8pt}

\begin{tcolorbox}[colback=soft,colframe=soft,arc=2mm,left=3mm,right=3mm,top=1.5mm,bottom=1.5mm]
\small Intenta resolver cada ejercicio antes de mirar su solución. Si no te sale el mismo resultado, recuerda que una derivada se puede escribir de muchas maneras equivalentes: compruébala con la calculadora de derivadas de los apuntes interactivos de la unidad.
\end{tcolorbox}

% =====================================================
\seccio{1}{Tasa de variación y derivada por definición}

\ex{1} \begin{tasks}(3)
  \task $\frac{6-0}{2}=\res{3}$
  \task $\frac{1-0}{e-1}=\res{\frac{1}{e-1}}$
  \task $\frac{1-0}{\pi/2}=\res{\frac2\pi}$
\end{tasks}

\ex{2} $v_m=\frac{s(2)-s(0)}{2}=\frac{20}{2}=\res{10\text{ m/s}}$. \ $\dfrac{s(1+h)-s(1)}{h}=\dfrac{-5(1+h)^2+20(1+h)-15}{h}=\dfrac{10h-5h^2}{h}=10-5h\to\res{10\text{ m/s}}$.

\ex{2} \begin{tasks}(1)
  \task $\dfrac{\frac{1}{(2+h)^2}-\frac14}{h}=\dfrac{4-(2+h)^2}{4h(2+h)^2}=\dfrac{-4-h}{4(2+h)^2}\to\res{-\frac14}$
  \task $\dfrac{\frac{1+h}{2+h}-\frac12}{h}=\dfrac{h}{2h(2+h)}\to\res{\frac14}$
  \task $\dfrac{e^{3h}-1}{h}\sim\dfrac{3h}{h}=\res{3}$
  \task $\dfrac{\sin\left(\frac\pi2+h\right)-1}{h}=\dfrac{\cos h-1}{h}\sim\dfrac{-h^2/2}{h}\to\res{0}$
\end{tasks}

\ex{3} \begin{tasks}(1)
  \task $\dfrac{(x+h)^3-x^3}{h}=3x^2+3xh+h^2\to\res{3x^2}$
  \task $\dfrac{\frac{1}{\sqrt{x+h}}-\frac1{\sqrt x}}{h}=\dfrac{\sqrt x-\sqrt{x+h}}{h\sqrt x\sqrt{x+h}}=\dfrac{-1}{\sqrt x\sqrt{x+h}\left(\sqrt x+\sqrt{x+h}\right)}\to\res{-\dfrac{1}{2x\sqrt x}}$
  \task $\dfrac{\cos x(\cos h-1)-\sin x\sin h}{h}\to\cos x\cdot0-\sin x\cdot1=\res{-\sin x}$
  \task $\dfrac{\ln(2x+2h)-\ln(2x)}{h}=\dfrac{\ln\left(1+\frac hx\right)}{h}\sim\dfrac{h/x}{h}=\res{\dfrac1x}$
\end{tasks}

\ex{2} \pau{Navarra, extraordinaria 2023}$g(1+h)=6+2h-h^2$, \ $g(1)=6$: \ $g'(1)=\displaystyle\lim_{h\to0}\frac{2h-h^2}{h}=\res{2}$


% =====================================================
\seccio{2}{Rectas tangente y normal}

\ex{1} \begin{tasks}(1)
  \task $f(2)=3$, $f'(x)=3x^2-3$, $f'(2)=9$: \ tangente $\res{y=9x-15}$; \ normal $y-3=-\frac19(x-2)\Rightarrow\res{y=-\frac19x+\frac{29}{9}}$
  \task $g(2)=\ln1=0$, $g'(x)=\frac{2x}{x^2-3}$, $g'(2)=4$: \ tangente $\res{y=4x-8}$; \ normal $\res{y=-\frac14x+\frac12}$
\end{tasks}

\ex{2} \begin{tasks}(1)
  \task $f(1)=\frac12$, $f'(x)=\frac{-1}{(x+1)^2}$, $f'(1)=-\frac14$: \ $\res{y=-\frac14x+\frac34}$
  \task $f(-2)=3$, $f'(x)=\frac{x}{\sqrt{5+x^2}}$, $f'(-2)=-\frac23$: \ $\res{y=-\frac23x+\frac53}$
  \task $f(1)=-\frac13$, $f'(x)=\frac{-(x^2+4)}{(x^2-4)^2}$, $f'(1)=-\frac59$: \ $\res{y=-\frac59x+\frac29}$
\end{tasks}

\ex{2} \pau{Islas Baleares, extraordinaria 2021}\begin{tasks}(1)
  \task Pendiente $-\frac34$: \ $-\frac{3}{x^2}=-\frac34\Rightarrow x=\pm2$: \ $\res{\left(2,\frac{19}{2}\right)}$ y $\res{\left(-2,\frac{13}{2}\right)}$
  \task $\res{y=-\frac34x+11}$ \ y \ $\res{y=-\frac34x+5}$
\end{tasks}

\ex{2} \pau{Murcia, septiembre 2018}$f'(x)=15x^2e^{2x}+10x^3e^{2x}-\dfrac{2x}{(x^2+1)^2}\Rightarrow\res{f'(0)=0}$: \ la tangente es horizontal, $\res{y=1}$.

\ex{2} $f'(x)=3x^2-12x+9=0\Rightarrow x=1,\ x=3$: \ $\res{(1,4)}$ y $\res{(3,0)}$. \ $3x^2-12x+9=-3\Rightarrow(x-2)^2=0$: en $(2,2)$, \ $\res{y=-3x+8}$.

\ex{3} Tangente en $(a,a^2)$: $y=2ax-a^2$. Pasa por $P$: $-3=2a-a^2\Rightarrow a=3$ o $a=-1$: \ $\res{y=6x-9}$ \ y \ $\res{y=-2x-1}$.

\ex{2} \begin{tasks}(1)
  \task $2a+\frac b2=4$, \ $a-\frac b4=0$: \ $\res{a=1,\ b=4}$
  \task $f(0)=b=1$; \ $f'(x)=\frac{a(x^2+1)-(ax+b)2x}{(x^2+1)^2}$, $f'(0)=a=1$: \ $\res{a=1,\ b=1}$
  \task $h'(1)=6+2a=0\Rightarrow a=-3$; \ $h(1)=1+a+3+b=2\Rightarrow\res{a=-3,\ b=1}$
\end{tasks}

\ex{2} \pau{La Rioja, ordinaria 2024}$y'=-\frac{1}{x^2}$, $y'(3)=-\frac19$: \ $\res{y=-\frac19x+\frac23}$. Corta los ejes en $(6,0)$ y $\left(0,\frac23\right)$, y el punto medio del segmento es $\left(\frac{6+0}{2},\frac{0+2/3}{2}\right)=\left(3,\frac13\right)$.

\ex{2} \pau{Castilla-La Mancha, extraordinaria 2021}$f(1)=\frac13$, $f'(x)=\frac{-6x^2+24x+6}{(3x^2+3)^2}\Rightarrow f'(1)=\frac23$: \ tangente $\res{y=\frac23x-\frac13}$; \ normal $\res{y=-\frac32x+\frac{11}{6}}$.

\ex{2} \pau{Andalucía, extraordinaria 2024}$f\left(\frac\pi2\right)=a+c=1$, \ $f'\left(\frac\pi2\right)=-b=0$, \ $f(0)=a+b=-1$: \ $\res{a=-1,\ b=0,\ c=2}$.

\ex{3} \pau{La Rioja, extraordinaria 2024}En $x=0$: $f'(0^-)=-1\neq1=f'(0^+)$, punto anguloso: \rest{no hay tangente}. En $x=-1$: $f(x)=-xe^{-x}$, $f(-1)=e$, $f'(-1)=-2e$: \ $\res{y=-2ex-e}$.

\ex{2} \pau{Aragón, ordinaria 2022}$f(-4)=2+e^3$, $f'(x)=\frac{-1}{2\sqrt{-x}}\Rightarrow f'(-4)=-\frac14$: \ $\res{y=-\frac14x+1+e^3}$.


% =====================================================
\seccio{3}{Reglas de derivación y derivadas sucesivas}

\ex{1} \begin{tasks}(2)
  \task $\left(x^{7/2}-2x^{-1/3}\right)^{\prime}=\res{\frac72x^2\sqrt x+\frac{2}{3x\sqrt[3]{x}}}$
  \task $\res{e^x(\sin x+x\sin x+x\cos x)}$
  \task $\res{\dfrac{4x}{(x^2+1)^2}}$
  \task $\res{\dfrac{x^2(2x-3)}{(x-1)^2}}$
  \task $\dfrac{\sin^2x-(1-\cos x)\cos x}{\sin^2x}=\dfrac{1-\cos x}{1-\cos^2x}=\res{\dfrac{1}{1+\cos x}}$
  \task $e^x(x^2-2x+2+2x-2)=\res{x^2e^x}$
\end{tasks}

\ex{1} \pau{Cantabria, ordinaria 2024}$f'(x)=1\cdot\ln x+x\cdot\frac1x=\res{\ln x+1}$

\ex{2} $F'(x)=\arctan x+\dfrac{x}{1+x^2}-\dfrac12\cdot\dfrac{2x}{x^2+1}=\res{\arctan x}$ \ (ya lo veremos: $F$ es una \emph{primitiva} de $\arctan x$).

\ex{2} \begin{tasks}(1)
  \task $f'=4x^3-6x^2+1$, $f''=12x^2-12x$, $\res{f'''=24x-12}$
  \task $\res{(-2)^ne^{-2x}}$; \ $\res{(x+n)e^x}$; \ $\res{\dfrac{(-1)^n\,n!}{(x+1)^{n+1}}}$; \ $\res{\dfrac{(-1)^{n+1}(n-1)!}{x^n}}$
  \task $2027=4\cdot506+3$: la misma que la tercera, $\res{\sin x}$
\end{tasks}

\ex{2} $f''=2a=4$, \ $f'(1)=2a+b=3$, \ $f(1)=a+b+c=2$: \ $\res{a=2,\ b=-1,\ c=1}$

\ex{3} $y'=e^{-x}(\cos x-\sin x)$, \ $y''=-2e^{-x}\cos x$. \ $y''+2y'+2y=e^{-x}\left(-2\cos x+2\cos x-2\sin x+2\sin x\right)=\res{0}$


% =====================================================
\seccio{4}{Regla de la cadena}

\ex{1} \begin{tasks}(3)
  \task $\res{24x(3x^2-1)^3}$
  \task $\res{\dfrac{x-1}{\sqrt{x^2-2x}}}$
  \task $\res{\dfrac{-6x}{(x^2+1)^4}}$
  \task $\res{\dfrac{2}{3\sqrt[3]{(2x+5)^2}}}$
  \task $\res{3e^{3x-1}}$
  \task $\res{2x\cdot2^{x^2}\ln2}$
\end{tasks}

\ex{2} \begin{tasks}(2)
  \task $2\sin3x\cos3x\cdot3=\res{3\sin6x}$
  \task $\res{-2x\sin(x^2)}$
  \task $\res{3\tan^2x\,(1+\tan^2x)}$
  \task $\res{\dfrac{2}{1+4x^2}}$
  \task $\dfrac{\frac{1}{2\sqrt x}}{\sqrt{1-x}}=\res{\dfrac{1}{2\sqrt x\sqrt{1-x}}}$
  \task $\dfrac{-(-e^{-x})}{\sqrt{1-e^{-2x}}}=\res{\dfrac{e^{-x}}{\sqrt{1-e^{-2x}}}}$
\end{tasks}

\ex{2} \begin{tasks}(2)
  \task $\res{\dfrac{2x-3}{x^2-3x}}$
  \task $\frac12\left[\frac{1}{1+x}+\frac{1}{1-x}\right]=\res{\dfrac{1}{1-x^2}}$
  \task $\res{\dfrac{3x^2}{(x^3+1)\ln2}}$
  \task $\res{-\tan x}$
  \task $\dfrac{1+\frac{x}{\sqrt{x^2+1}}}{x+\sqrt{x^2+1}}=\res{\dfrac{1}{\sqrt{x^2+1}}}$
  \task $\dfrac{\frac12\left(1+\tan^2\frac x2\right)}{\tan\frac x2}=\dfrac{1}{2\sin\frac x2\cos\frac x2}=\res{\dfrac{1}{\sin x}}$
\end{tasks}

\ex{2} \begin{tasks}(2)
  \task $\res{(1-2x^2)\,e^{-x^2}}$
  \task $\res{(x-1)(x+2)^2(5x+1)}$
  \task $\res{\dfrac{2\sqrt x+1}{4\sqrt x\,\sqrt{x+\sqrt x}}}$
  \task $\res{e^{\sin x}\left(\cos^2x-\sin x\right)}$
  \task $\res{\dfrac{\cos(\ln x)-\sin(\ln x)}{x}}$
\end{tasks}

\ex{2} \pau{Murcia, septiembre 2020}\begin{tasks}(2)
  \task $\res{2x\ln x+\dfrac{x^2-2}{x}}$
  \task $\res{12x^2e^{4x^3+2}}$
\end{tasks}

\ex{2} \pau{Murcia, junio 2018}\begin{tasks}(2)
  \task $\res{2x\ln(x^2)+2x}$
  \task $\res{(2x-3)\,e^{-3x+x^2}}$
\end{tasks}

\ex{2} \pau{Andalucía, ordinaria 2021}\begin{tasks}(2)
  \task $\frac{1}{x-1}-\frac{1}{x+1}=\res{\dfrac{2}{x^2-1}}$
  \task $\res{e^{2x^2}\left(3x^2+4x^4\right)}$
\end{tasks}

\ex{2} \pau{Andalucía, ordinaria (suplente) 2025}\begin{tasks}(2)
  \task $\res{\dfrac{4x}{(x^2-1)(x^2+1)}}$
  \task $\res{(4x^2-4x+3)\,e^{x^2-x}}$
\end{tasks}

\ex{3} \pau{Andalucía, ordinaria 2024}\begin{tasks}(1)
  \task $\res{-2(x^2-3x+2)(x^2+2)^2e^{-2x}}$
  \task $\res{\dfrac{-3x^2(1-2x^2)+8x(1-x^3)\ln(1-x^3)}{(1-x^3)(1-2x^2)^3}}$
\end{tasks}

\ex{2} \pau{Murcia, septiembre 2018}$f(0)=b=2$; \ $f'(x)=3x^2\ln(2x+5)+\dfrac{2x^3}{2x+5}+a$, \ $f'(0)=a=1$: \ $\res{a=1,\ b=2}$


% =====================================================
\seccio{5}{Derivación logarítmica}

\ex{2} \begin{tasks}(1)
  \task $\ln y=x\ln x\Rightarrow\res{y^{\prime}=x^x(\ln x+1)}$
  \task $\res{y^{\prime}=x^{\sin x}\left(\cos x\ln x+\dfrac{\sin x}{x}\right)}$
  \task $\ln y=(\ln x)^2\Rightarrow\res{y^{\prime}=\dfrac{2\ln x}{x}\,x^{\ln x}}$
  \task $\ln y=\dfrac{\ln x}{x}\Rightarrow\res{y^{\prime}=x^{1/x}\,\dfrac{1-\ln x}{x^2}}$
  \task $\res{y^{\prime}=(x^2+1)^x\left(\ln(x^2+1)+\dfrac{2x^2}{x^2+1}\right)}$
  \task $\res{y^{\prime}=(\ln x)^x\left(\ln(\ln x)+\dfrac{1}{\ln x}\right)}$
  \task $\res{y^{\prime}=(\cos x)^{\sin x}\left(\cos x\ln(\cos x)-\dfrac{\sin^2x}{\cos x}\right)}$
  \task $\ln y=x^2\ln x\Rightarrow\res{y^{\prime}=x^{x^2+1}(2\ln x+1)}$
\end{tasks}

\ex{2} \begin{tasks}(1)
  \task $\res{y^{\prime}=\dfrac{x^3(x+1)^5}{\sqrt{x^2+4}}\left(\dfrac3x+\dfrac{5}{x+1}-\dfrac{x}{x^2+4}\right)}$
  \task $\res{y^{\prime}=\dfrac{(x-1)^2\sqrt[3]{x+2}}{(x+3)^4}\left(\dfrac{2}{x-1}+\dfrac{1}{3(x+2)}-\dfrac{4}{x+3}\right)}$
\end{tasks}

\ex{2} \begin{tasks}(1)
  \task $f'(x)=x^x(\ln x+1)$: \ $f(1)=1$, $f'(1)=1$: $\res{y=x}$; \ $f(2)=4$, $f'(2)=4(1+\ln2)$: $\res{y=4(1+\ln2)(x-2)+4}$
  \task $\ln x+1=0\Rightarrow x=\frac1e$: \ $\res{\left(\frac1e,\ e^{-1/e}\right)}$
  \task La regla de la potencia exige exponente constante, y la de la exponencial, base constante. Aquí varían ambas: de hecho, $(x^x)'$ es la \rest{suma} de las dos expresiones incorrectas, $x\cdot x^{x-1}+x^x\ln x$.
\end{tasks}


% =====================================================
\seccio{6}{Derivación implícita}

\ex{2} \begin{tasks}(2)
  \task $2x+2yy^{\prime}=4\Rightarrow\res{y^{\prime}=\dfrac{2-x}{y}}$
  \task $\res{y^{\prime}=\dfrac{y-x^2}{y^2-x}}$
  \task $\res{y^{\prime}=-\dfrac{2xy+y^2}{x^2+2xy}}$
  \task $e^{xy}(y+xy^{\prime})=1+y^{\prime}\Rightarrow\res{y^{\prime}=\dfrac{1-ye^{xy}}{xe^{xy}-1}}$
  \task $\res{y^{\prime}=-\dfrac{\sqrt y}{\sqrt x}}$
  \task $1=y^{\prime}+\dfrac{y^{\prime}}{y}\Rightarrow\res{y^{\prime}=\dfrac{y}{y+1}}$
\end{tasks}

\ex{2} \begin{tasks}(1)
  \task $y^{\prime}=-\frac{2x+y}{x+2y}=-\frac54$: \ tangente $\res{y=-\frac54x+\frac72}$; \ normal $\res{y=\frac45x-\frac35}$
  \task $y^{\prime}=\frac{3y-x^2}{y^2-3x}=\frac{8}{10}$: \ $\res{y=\frac45x+\frac{12}{5}}$
  \task $9x^2+16y^2=144$: $y^{\prime}=-\frac{9x}{16y}=-1$: \ $\res{y=-x+5}$
  \task $2yy^{\prime}=3x^2$: $y^{\prime}=\frac32$: \ $\res{y=\frac32x-\frac12}$
\end{tasks}

\ex{3} $y'=-\dfrac{2x+y}{x+2y}$. Horizontal si $2x+y=0$: $3x^2=3$, \ $\res{(1,-2)}$ y $\res{(-1,2)}$. Vertical si $x+2y=0$ (denominador nulo): $3y^2=3$, \ $\res{(-2,1)}$ y $\res{(2,-1)}$.

\ex{3} \begin{tasks}(1)
  \task $y^{\prime}=-\frac xy$, \ $y^{\prime\prime}=-\frac{y-xy^{\prime}}{y^2}=\res{-\dfrac{r^2}{y^3}}$
  \task $y^{\prime}=\frac xy$, \ $\res{y^{\prime\prime}=\dfrac{y^2-x^2}{y^3}=-\dfrac{1}{y^3}}$
  \task $y^{\prime}=\frac{1}{3y^2+1}$, \ $\res{y^{\prime\prime}=\dfrac{-6y\,y^{\prime}}{(3y^2+1)^2}=\dfrac{-6y}{(3y^2+1)^3}}$
\end{tasks}

\ex{2} \begin{tasks}(1)
  \task $e^yy^{\prime}=1\Rightarrow\res{y^{\prime}=\frac{1}{e^y}=\frac1x}$
  \task $\cos y\,y^{\prime}=1\Rightarrow\res{y^{\prime}=\frac{1}{\sqrt{1-\sin^2y}}=\frac{1}{\sqrt{1-x^2}}}$
  \task $(1+\tan^2y)\,y^{\prime}=1\Rightarrow\res{y^{\prime}=\frac{1}{1+x^2}}$
\end{tasks}

\ex{3} $y^{\prime}\ln x+\dfrac yx=\ln y+\dfrac{x\,y^{\prime}}{y}$. En $(2,4)$: $y^{\prime}\left(\ln2-\frac12\right)=2\ln2-2$: \ $\res{y^{\prime}=\dfrac{4(\ln2-1)}{2\ln2-1}\approx-3{,}18}$


% =====================================================
\seccio{7}{Continuidad y derivabilidad}

\ex{2} \begin{tasks}(2)
  \task $-1\neq1$: \rest{no} (anguloso)
  \task $-6\neq6$: \rest{no}
  \task $-1\neq1$: \rest{no}
  \task $\frac{h^{2/3}}{h}\to\mp\infty$: \rest{no} (retroceso)
  \task $-1\neq1$: \rest{no}
  \task $\frac{h^2|h|}{h}\to0$: \rest{sí}, $f'(0)=0$
\end{tasks}

\ex{1} \begin{tasks}(1)
  \task \rest{Falso}: $|x|$ en $x=0$.
  \task \rest{Verdadero}: es el contrarrecíproco de «derivable $\Rightarrow$ continua».
  \task \rest{Falso}: $|x|+(-|x|)=0$.
  \task \rest{Falso}: $f(x)=x$ es derivable en 0 y $|x|$ no.
\end{tasks}

\ex{2} \pau{Galicia, ordinaria 2026}\rest{Es posible}: por ejemplo $f(x)=|x|$ (punto anguloso en el origen) o $f(x)=\sqrt[3]{x}$ (tangente vertical).

\ex{3} $\dfrac{f(h)-f(0)}{h}=h\sin\dfrac1h\to0$ (0 por acotada): \ $\res{f'(0)=0}$. \ En cambio, $\dfrac{g(h)}{h}=\sin\dfrac1h$ oscila entre $-1$ y $1$: \ $g$ \rest{no es derivable} en 0 (aunque es continua).


% =====================================================
\seccio{8}{Derivabilidad de funciones a trozos}

\ex{2} \begin{tasks}(2)
  \task $1=1$; $f'(0^-)=e^0=1=f'(0^+)$: \rest{derivable en $\mathbb R$}
  \task $2=2$; $g'(1^-)=-2=-\frac{2}{1^2}=g'(1^+)$: \rest{derivable en $\mathbb R$}
\end{tasks}

\ex{2} \pau{Andalucía, extraordinaria 2024}$x=4$: $3=3$, continua. $f'(4^-)=-2\cdot4+4=-4\neq2=f'(4^+)$: \ continua en $\mathbb R$, \rest{derivable en $\mathbb R-\{4\}$}.

\ex{2} \pau{Andalucía, ordinaria (suplente) 2025}$f(-2)=5=\displaystyle\lim_{x\to-2^+}(x^2+1)$: continua. $f'(-2^-)=\frac52\neq-4=f'(-2^+)$: \ \rest{no derivable en $x=-2$}.

\ex{2} \pau{Andalucía, ordinaria 2024}$t=10$: $80=80$. \ $v'(10^-)=2\cdot10-8=12$, \ $v'(10^+)=-20+32=12$: \ \rest{continua en $[0,24]$ y derivable en $(0,24)$}, con $v'(10)=12$.

\ex{2} \pau{Navarra, ordinaria 2024}$x=-1$: $3\neq5$, \rest{salto finito} (no derivable). $x=1$: $3=3$, continua; $f'(1^-)=-1\neq-4=f'(1^+)$. \ \rest{Continua en $\mathbb R-\{-1\}$, derivable en $\mathbb R-\{-1,1\}$}.

\ex{2} \pau{Andalucía, extraordinaria 2021}$\displaystyle\lim_{x\to0^-}2^{x+1}=2\neq0=f(0)$: salto finito en $x=0$. \ \rest{Continua y derivable en $\mathbb R-\{0\}$}.

\ex{3} \pau{Madrid, extraordinaria 2022 (ampliado)}$\displaystyle\lim_{x\to0^+}x\ln x=0=f(0)$: \rest{continua}. $f'(0^-)=1$, \ $f'(0^+)=\displaystyle\lim_{x\to0^+}(\ln x+1)=-\infty$: \ \rest{no derivable en $x=0$}.

\ex{2} \pau{Asturias, ordinaria 2025 (ampliado)}$x=4$: $2=2$, continua en $[0,+\infty)$. $f'(4^-)=\frac12\neq2=f'(4^+)$: \ \rest{derivable en $(0,4)\cup(4,+\infty)$}.

\ex{3} \pau{Andalucía, ordinaria 2022 (ampliado)}\begin{tasks}(1)
  \task $-1=-a+b$ y $a+b=\frac12$: \ $\res{a=\frac34,\ b=-\frac14}$
  \task $x=-1$: $f'(-1^-)=-\frac{1}{(-1)^2}=-1\neq\frac34$; \ $x=1$: $f'(1^-)=\frac34$ y $f'(1^+)=\frac{x^2+2x}{(x+1)^2}\big|_{x=1}=\frac34$. \ \rest{Derivable en $\mathbb R-\{-1\}$}.
\end{tasks}

\ex{2} \begin{tasks}(2)
  \task $0=0$; $f'(0^\pm)=1$: \rest{derivable en $(-1,+\infty)$}
  \task $1=1$; $g'(0^-)=-\sin0=0=g'(0^+)$: \rest{derivable en $\mathbb R$}
\end{tasks}

\ex{3} \pau{Andalucía, extraordinaria 2021}$t=1{,}8$: $0{,}06=0{,}06$; $f'(1{,}8^-)=-1{,}6\neq0{,}1$. \ $t=5$: $0{,}38=0{,}38$; $f'(5^-)=0{,}1\neq3{,}3$. \ \rest{Continua en $[0{,}2;\,10]$, derivable en $(0{,}2;\,10)-\{1{,}8;\ 5\}$}.

\ex{3} \begin{tasks}(1)
  \task $f=-2x$ ($x<-1$), $2$ ($-1\leq x\leq1$), $2x$ ($x>1$): \rest{derivable en $\mathbb R-\{-1,1\}$}
  \task $g'(2^-)=-2\neq2=g'(2^+)$: \rest{$\mathbb R-\{2\}$}
  \task $h'(0^-)=1\neq-1=h'(0^+)$: \rest{$\mathbb R-\{0\}$}
  \task $x^2-5x+4<0$ en $(1,4)$; $k'(1^\mp)=\mp3$, $k'(4^\mp)=\mp3$: \rest{$\mathbb R-\{1,4\}$}
\end{tasks}


% =====================================================
\seccio{9}{Funciones a trozos con parámetros}

\ex{2} \pau{Andalucía, ordinaria 2022}$3=a+b+2$ \ y \ $6=2a+b$: \ $\res{a=5,\ b=-4}$

\ex{2} \pau{Andalucía, extraordinaria 2020}$2-a=a+b$; \ $f'(0^-)=-a$, $f'(0^+)=b$: \ $\res{a=2,\ b=-2}$

\ex{2} \begin{tasks}(1)
  \task Continua para todo $a$; $-3=a$: \ $\res{a=-3}$
  \task Continua para todo $a$ ($2=e^0+1$); $-3=a\,e^0$: \ $\res{a=-3}$
\end{tasks}

\ex{2} \pau{Islas Baleares, extraordinaria 2021}\begin{tasks}(1)
  \task $x=-2$: $8+a=-1\Rightarrow\res{a=-9}$ ($x=1$: $-4=-4$ siempre)
  \task Para $a=1$ no es continua en $x=-2$ $\Rightarrow$ no derivable. En $x=1$: $f'(1^-)=2\neq-7$. \ \rest{No es derivable ni en $-2$ ni en $1$}.
\end{tasks}

\ex{3} \pau{Aragón, ordinaria 2022}\begin{tasks}(1)
  \task $2\sqrt3=\sqrt b$ y $\sqrt{3a+b}=3$: \ $\res{a=-1,\ b=12}$
  \task $f'(0^-)=0\neq-\frac{1}{4\sqrt3}=f'(0^+)$: \rest{no derivable en 0}; \ $f'(3^\pm)=-\frac16$: \rest{derivable en 3}
\end{tasks}

\ex{2} \pau{Andalucía, ordinaria (suplente) 2025}$\res{f(0)=10}$. \ $8{,}75=6{,}25+2{,}5a+b$ \ y \ $-3=5+a$: \ $\res{a=-8,\ b=22{,}5}$

\ex{3} \pau{Cataluña, ordinaria 2026 (ampliado)}$x=0$: $m=\pm2$; \ $x=2$: $m=-2$: \ $\res{m=-2}$. \ $f'(0^-)=10\neq-4=f'(0^+)$; \ $f'(2^\pm)=0$: \ \rest{derivable en $\mathbb R-\{0\}$}.

\ex{2} \pau{Extremadura, extraordinaria 2023 (ampliado)}$2+a+b=0$, \ $2-a+b=5$: \ $\res{a=-\frac52,\ b=\frac12}$. \ $f'(1^-)=4+a=\frac32\neq1=f'(1^+)$: \ \rest{no derivable en $x=1$}.

\ex{2} \pau{Aragón, extraordinaria 2021 (ampliado)}$5-a=\frac6a\Rightarrow\res{a=2\ \text{o}\ a=3}$. \ $f'(1^-)=-2a$, $f'(1^+)=-\frac6a$: \ $-4\neq-3$, \ $-6\neq-2$: \ \rest{nunca es derivable en $x=1$}.

\ex{2} \begin{tasks}(2)
  \task $a+b=0$, $2a+b=1$: $\res{a=1,\ b=-1}$
  \task $a+b=1$, $2a=-1$: $\res{a=-\frac12,\ b=\frac32}$
  \task $1+b=3$, $a=2$: $\res{a=2,\ b=2}$
  \task $a=0$, $b=\frac{1}{1+0}$: $\res{a=0,\ b=1}$
\end{tasks}

\ex{2} \begin{tasks}(1)
  \task Continua si $k+2=1-m$, es decir, $\res{m=-k-1}$. Derivable si, además, $2k+2=2$: $\res{k=0,\ m=-1}$.
  \task $f(0)=-1$; el límite por la derecha es finito solo si $k=0$, y entonces $\lim_{x\to0^+}(-e^x)=-1$: $\res{k=0}$. $f'(0^-)=b$, $f'(0^+)=-1$: $\res{b=-1,\ k=0}$.
  \task $a-1=b$, $\ -1=-b-1$: $\res{a=1,\ b=0}$; \ $f(x)=1-x$ si $x\leq1$ y $-\ln x$ si $x>1$.
\end{tasks}

\ex{2} \begin{tasks}(1)
  \task $b=1$, $\ -a=b$: $\res{a=-1,\ b=1}$
  \task $\frac{x^2-1}{x-1}=x+1$: continuidad $a+b=2$, derivabilidad $a=1$: $\res{a=b=1}$
  \task Continuidad: $3-k=\frac2k\Rightarrow k=1$ o $k=2$; derivabilidad: $-2k=-\frac2k\Rightarrow k=\pm1$: $\res{k=1}$
\end{tasks}

\ex{3} \pau{Madrid, ordinaria 2026}\begin{tasks}(1)
  \task Continuidad: $b=2$. $f'(0^-)=-a$, $f'(0^+)=\frac{2e^0(1+0)-2e^0\cdot0}{1}=2$: $\res{a=-2,\ b=2}$.
  \task $f(2)=\frac{2e^2}{5}$, $f'(x)=\frac{2e^x(x-1)^2}{(1+x^2)^2}\Rightarrow f'(2)=\frac{2e^2}{25}$: $\res{y=\frac{2e^2}{25}x+\frac{6e^2}{25}}$.
\end{tasks}


\end{document}
